\SourcePage{246}{249}
\section{Spin}

In classical and quantum mechanics alike, conservation of angular momentum
follows from the isotropy of space for a closed system. This already displays
the connection between angular momentum and rotational symmetry. In quantum
mechanics, however, that connection becomes especially profound and is in
essence the central content of the concept of angular momentum. Moreover, the
classical definition of a particle's angular momentum as the product
$\mathbf r\times\mathbf p$ loses its direct meaning because position and
momentum cannot be measured simultaneously.

We saw in \S~28 that specifying $l$ and $m$ fixes the angular dependence of a
particle's wave function, and thereby all its symmetry properties under
rotations. In their most general formulation, these properties amount to the
transformation law of wave functions when the coordinate system is rotated.

The wave function $\psi_{LM}$ of a system of particles, with specified angular
momentum $L$ and projection $M$, remains unchanged\footnote{Apart from an
irrelevant phase factor.} only when the coordinate system is rotated about the
$z$ axis. Any rotation that changes the direction of that axis makes the
angular-momentum projection on it indefinite. In the new coordinate axes the
wave function therefore becomes, in general, a superposition (linear
combination) of the $2L+1$ functions belonging to the possible values of $M$
at fixed $L$. Thus, under rotations, the $2L+1$ functions $\psi_{LM}$
transform among themselves\footnote{In mathematical terminology, these
functions realize the so-called \emph{irreducible representations} of the
rotation group. The number of functions transforming among themselves is
called the \emph{dimension of the representation}; it is understood that no
choice of other linear combinations can reduce this number.}. The law of
this transformation, that is, the superposition coefficients as functions of
the angles through which the coordi\SourcePageJoin{247}{250}nate axes are
rotated, is completely fixed by $L$. Angular momentum thus acquires the
meaning of a quantum number that classifies the states of a system by their
transformation properties under coordinate rotations. This aspect is
particularly important in quantum mechanics because it is not tied directly
to an explicit angular dependence of the wave functions: their mutual
transformation law can be stated without reference to such a dependence.

Consider a composite particle, say an atomic nucleus, at rest as a whole and
in a definite internal state. Besides a definite internal energy it has an
angular momentum of definite magnitude $L$, associated with the motion of its
constituents; this angular momentum may have $2L+1$ different orientations in
space. Hence, in describing the motion of the composite particle as a whole,
we must assign to it, in addition to its coordinates, a discrete variable:
the projection of its internal angular momentum on a chosen direction.

With the meaning of angular momentum just described, its origin is immaterial.
We are led naturally to the idea of an ``intrinsic'' angular momentum that
must be assigned to a particle regardless of whether the particle is
``composite'' or ``elementary.''

Quantum mechanics therefore attributes to an elementary particle an
``intrinsic'' angular momentum unrelated to its motion through space. This is
a specifically quantum property of elementary particles, disappearing in the
limit $\hbar\to0$, and in principle permits no classical interpretation
\footnote{In particular, it would be wholly meaningless to imagine the
``intrinsic'' angular momentum of an elementary particle as arising from its
rotation ``about its own axis.''}.

The intrinsic angular momentum of a particle is called its \emph{spin}, in
contrast to the angular momentum associated with motion through space, which
is called \emph{orbital angular momentum}\footnote{The physical idea that the
electron possesses intrinsic angular momentum was proposed by G.~Uhlenbeck
and S.~Goudsmit (1925). Spin was introduced into quantum mechanics by
W.~Pauli (1927).}. The particle may be elementary, or it may be composite but
behave as elementary in the class of phenomena under consideration (an atomic
nucleus, for example). We denote its spin, measured like orbital angular
momentum in units of $\hbar$, by $s$.

\SourcePage{248}{251}
For a particle with spin, the wave-function description of a state must give
not only the probabilities of its various positions in space but also the
probabilities of the possible orientations of its spin. Thus the wave function
depends not only on the three continuous particle coordinates, but also on a
discrete \emph{spin variable}. This variable specifies the spin projection on
a chosen direction (the $z$ axis) and takes a finite set of discrete values,
which will be denoted by $\sigma$.

Let $\psi(x,y,z;\sigma)$ be such a wave function. It is essentially a set of
different coordinate functions, one for each value of $\sigma$; we shall call
them the \emph{spin components} of the wave function. The integral
\[
 \int |\psi(x,y,z;\sigma)|^2\,dV
\]
is the probability that the particle has the specified value $\sigma$. The
probability that it lies in the volume element $dV$ with any value of
$\sigma$ is
\[
 dV\sum_\sigma |\psi(x,y,z;\sigma)|^2.
\]

When applied to a wave function, the quantum-mechanical spin operator acts on
the spin variable $\sigma$. In other words, it transforms the wave-function
components among themselves in a certain way. Its form will be obtained
below. Even from general considerations, however, it is easy to see that
$\hat s_x$, $\hat s_y$, and $\hat s_z$ obey the same commutation relations as
the orbital-angular-momentum operators.

The angular-momentum operator is essentially the operator of an infinitesimal
rotation. In deriving the orbital-angular-momentum operator in \S~26, we
considered how a rotation acts on a function of the coordinates. That
derivation is inapplicable to spin, since the spin operator acts on the spin
variable rather than on the coordinates. We must instead consider an
infinitesimal rotation in its general form, as a rotation of the coordinate
system. Carrying out infinitesimal rotations successively about the $x$ and
$y$ axes and then about the same axes in the reverse order, direct calculation
shows that the difference between the two results is equivalent to an
infinitesimal \SourcePage{249}{252}rotation about the $z$ axis, through an
angle equal to the product of the rotation angles about $x$ and $y$. We omit
this simple calculation. It again gives the usual commutation relations for
the components of angular momentum, and these must therefore hold for spin:
\begin{equation}
 \{\hat s_x,\hat s_y\}=i\hat s_z
 \tag{54.1}
\end{equation}
with all the physical consequences that follow from them.

The commutation relations (54.1) determine the possible magnitudes and
components of spin. The whole derivation in \S~27, equations (27.7)--(27.9),
used only the commutation relations and applies here without change, with $s$
in place of $L$. Equations (27.7) show that the eigenvalues of a spin
projection form a sequence differing by unity. Here, however, we cannot infer
that the values themselves must be integral, as for the orbital projection
$L_z$: the argument at the beginning of \S~27 used the specifically orbital
expression (26.14) for $\hat l_z$ and does not apply.

The sequence of eigenvalues of $s_z$ is bounded above and below by values of
equal magnitude and opposite sign, denoted by $\pm s$. The difference $2s$
between the greatest and least values must be an integer or zero. Hence $s$
may be $0$, $1/2$, $1$, $3/2$,~\dots

The eigenvalues of the square of the spin are consequently
\begin{equation}
 s^2=s(s+1),
 \tag{54.2}
\end{equation}
where $s$ may be integral, including zero, or half-integral. For given $s$,
the spin component $s_z$ can take the $2s+1$ values
$s,s-1,\ldots,-s$. Accordingly, the wave function of a spin-$s$ particle has
$2s+1$ components\footnote{For each kind of particle, $s$ is fixed. Thus in
the classical limit $\hbar\to0$, the spin angular momentum $\hbar s$ tends
to zero. This reasoning is inapplicable to orbital angular momentum because
$l$ can take arbitrary values. The classical limit involves simultaneously
$\hbar\to0$ and $l\to\infty$, with $\hbar l$ remaining finite.}.

\SourcePage{250}{253}
Experiment shows that most elementary particles---electrons, positrons,
protons, neutrons, $\mu$ mesons, and all hyperons ($\Lambda$, $\Sigma$,
$\Xi$)---have spin $1/2$. There are also elementary particles, the $\pi$ and
$K$ mesons, with spin zero.

The total angular momentum of a particle is the sum of its orbital angular
momentum $\mathbf l$ and spin $\mathbf s$. Since their operators act on
functions of wholly different variables, they of course commute.

The eigenvalues of the total angular momentum
\begin{equation}
 \mathbf j=\mathbf l+\mathbf s
 \tag{54.3}
\end{equation}
are determined by the same ``vector-model'' rule as the sum of the orbital
angular momenta of two different particles (\S~31). For fixed $l$ and $s$,
the total angular momentum can take the values
$l+s,l+s-1,\ldots,|l-s|$. Thus an electron of spin $1/2$ with nonzero
orbital angular momentum $l$ may have $j=l\pm1/2$; when $l=0$, naturally,
there is only $j=1/2$.

The total-angular-momentum operator $\mathbf J$ of a system of particles is
the sum of their individual operators $\mathbf j$, so its values are again
given by the vector-model rules. It may be written
\begin{equation}
 \mathbf J=\sum_a\mathbf j_a=\mathbf L+\mathbf S,
 \qquad \mathbf L=\sum_a\mathbf l_a,\quad
 \mathbf S=\sum_a\mathbf s_a,
 \tag{54.4}
\end{equation}
where $\mathbf S$ may be called the total spin and $\mathbf L$ the total
orbital angular momentum of the system.

If the total spin of the system is half-integral (or integral), the total
angular momentum is likewise half-integral (or integral), since orbital
angular momentum is always integral. In particular, a system containing an
even number of identical particles necessarily has integral total spin and
hence integral total angular momentum.

The total-angular-momentum operators $\mathbf j$ of a particle (or $\mathbf
J$ of a system) obey the same commutation rules as orbital angular momentum or
spin, since these are the general commutation rules for every angular
momentum. Equations (27.13) for its matrix elements likewise hold for every
angular momentum when the elements are defined with respect to its own
eigenfunctions. Equations (29.7)--(29.10) for matrix elements of arbitrary
vector quantities also remain valid, with the corresponding changes of
notation.

\SourcePage{251}{254}
\ProblemHeading

A spin-$1/2$ particle is in a state with the definite value $s_z=1/2$. Find
the probabilities of the possible values of its spin projection on an axis
$z'$ inclined at an angle $\theta$ to the $z$ axis.

\SolutionHeading

The mean spin vector $\bar{\mathbf s}$ plainly points along $z$ and has
magnitude $1/2$. Its projection on $z'$ gives the mean spin in that direction,
$\bar s_{z'}=(1/2)\cos\theta$. On the other hand,
$\bar s_{z'}=(1/2)(w_+-w_-)$, where $w_\pm$ are the probabilities of
$s_{z'}=\pm1/2$. Together with $w_++w_-=1$, this gives
\[
 w_+=\cos^2(\theta/2),\qquad w_-=\sin^2(\theta/2).
\]
